% TOP_PROBLEM: {"id":30007757,"problem_number":"LOCAL-30007757","title":"Global minimum taxes and real investment","category_id":21,"category":{"id":21,"name":"economics","display_name":"Economics","description":"Open economic theory statements, published research agendas, and proposed scoped specifications; question provenance and historical priority are qualified per record.","slug":"economics","order_index":21,"created_at":"2026-10-08T00:00:00Z"},"status":"provenance_pending","external_url":null,"legacy_ids":[],"published":false,"statement":"In the explicitly specified setting: How do global minimum taxes change profit shifting, investment locations, research activity, and countries' strategic tax responses? The mathematical statement above fixes the target, assumptions and scope of this catalogue formulation.","economics":{"original_id":246,"field":"Taxation and social insurance","kind":"design","formalization_basis":"proposed_specification","source_status":"not_verified","priority_status":"not_established","priority_note":"An original explicit open-question statement was not verified; the cited supporting paper must not be treated as its first listing.","earliest_verified_reference":null,"current_reference":{"authors":["Sebastian Dyrda","Guangbin Hong","Muhammad Ali Sajid","Joseph B. Steinberg"],"title":"Global Ripple Effects of Corporate Tax Reforms","year":2026,"url":"https://www.nber.org/papers/w34627","doi":"10.3386/w34627","locator":"Working Paper 34627; bibliographic record and abstract","evidence_summary":"Models international corporate tax reforms and spillovers. No explicit statement of the complete catalogue question as an unresolved problem was verified."},"formalization_note":"Catalogue-authored scoped formalization. Population, instruments, welfare objective and identification restrictions must be instantiated before claiming a resolution; source authors are not credited with these added assumptions."},"statusQualification":"Provenance pending: an explicit published statement of this exact open question was not verified. This is a catalogue-authored candidate specification, not a certified open conjecture. Catalogue-authored scoped formalization. Population, instruments, welfare objective and identification restrictions must be instantiated before claiming a resolution; source authors are not credited with these added assumptions.","statusReviewedAt":"2026-10-08"}
% FIRST_POSTED: 2026-10-08
% ENTRY_KIND: statement_only
% !TeX program = lualatex
\documentclass[11pt]{article}
\usepackage{fontspec}
\usepackage[a4paper,margin=25mm]{geometry}
\usepackage{amsmath,amssymb,mathtools}
\usepackage{xurl}
\usepackage[hidelinks]{hyperref}
\newcommand{\sref}[2]{\hyperlink{source:#1:#2}{[S#2]}}
\setlength{\emergencystretch}{3em}
\begin{document}
% problemId:30007757
\section*{Global minimum taxes and real investment}\label{30007757}
\subsection{Definitions and mathematical statement}
Consider \(J\) jurisdictions and heterogeneous multinational firms indexed by \(f\). A minimum-tax regime \(m\) specifies the effective-rate floor, top-up priority, covered firms, tax base and deductions. Governments choose domestic instruments \(t_j\) from specified compact sets. Firms choose tangible investment \(K_{fj}\), research spending \(D_{fj}\), and declared profits \(\Pi_{fj}\), subject to technologies, real-resource costs of shifting, enforcement and the exact legal accounting rules. Let \(\mathcal E(m)\) be equilibria in which firms optimize and each government maximizes a specified resident-welfare objective given other governments' instruments. Baseline \(m_0\) has no minimum. Specify an equilibrium-selection rule for both regimes, or report sets across their nonempty equilibrium correspondences. Determine the counterfactual vector
\[\Delta(m,m_0)=\big(\Delta\sum_f K_{fj},\ \Delta\sum_f D_{fj},\ \Delta\sum_f\Pi_{fj},\ \Delta R_j,\ \Delta W_j\big)_{j=1}^J,\]
Here \(\Delta X=X(m)-X(m_0)\) uses the selected equilibria, \(R_j\) is net revenue, and \(W_j\) resident welfare. Characterize parameter and implementation conditions under which lower profit shifting raises aggregate real investment or instead relocates it. Identify or bound primitives from firm ownership, production and tax records, and test predictions against implementation episodes. Distinguish fixed-domestic-rate responses from strategic government responses. This is a proposed model-and-identification target; solving an existing calibrated reform model does not establish general effects or original open-question priority.

\par\textbf{Formulation and scope.} Catalogue-authored scoped formalization. Population, instruments, welfare objective and identification restrictions must be instantiated before claiming a resolution; source authors are not credited with these added assumptions.

\par\textbf{Open-question provenance.} An explicit published open-question statement was not verified. Sources below are supporting literature and must not be treated as a first listing.

\par\textbf{Historical priority.} An original explicit open-question statement was not verified; the cited supporting paper must not be treated as its first listing.

\subsection{Short English statement}
In the explicitly specified setting: How do global minimum taxes change profit shifting, investment locations, research activity, and countries' strategic tax responses? The mathematical statement above fixes the target, assumptions and scope of this catalogue formulation.

\subsection{Sources}
\begin{itemize}\raggedright
\item[\textnormal{[S1]}]\hypertarget{source:30007757:1}{} Sebastian Dyrda, Guangbin Hong, Muhammad Ali Sajid, Joseph B. Steinberg. 2026. Global Ripple Effects of Corporate Tax Reforms. Working Paper 34627; bibliographic record and abstract.
\url{https://www.nber.org/papers/w34627}.
DOI: 10.3386/w34627.
{\small\textit{Supporting or current literature; see evidence qualification. Models international corporate tax reforms and spillovers. No explicit statement of the complete catalogue question as an unresolved problem was verified.}}
\end{itemize}
\end{document}
